\section{Four Evaluation Traps}
\label{sec:traps}

Our anatomy repeatedly reached conclusions opposite to those suggested by the metrics most common in this literature. We document four failure modes; each reversed at least one finding in this paper.

\paragraph{Trap 1: cosine and probe metrics are duals of the training loss.}
Probe correlation and latent cosine rise mechanically when probe or slot losses are optimized---they measure whether information is \emph{present}, not whether prediction \emph{uses} it. Probe supervision lifts velocity decodability from $\rho{=}0.44$ to $0.80$ and gives the most stable long-horizon \emph{cosine} ($0.882$ vs.\ $0.843$), while its \emph{pixel} rollout is the worst in the comparison ($19.93$ vs.\ $20.64$\,dB at horizon $28$). Appearance augmentation on Physion++ improves the aggregate long-horizon cosine while per-scene nMSE collapses up to $100\times$. Judging either result by direction-only or readout metrics inverts the conclusion.

\paragraph{Trap 2: nMSE denominators can degenerate.}
The scale-aware metric has its own failure: when the target sequence degenerates (a parabola ball exits the frame, target variance ${\to}0$), per-horizon nMSE explodes---up to $2{\times}10^{6}$ at horizon $28$ while cosine sits at a healthy $0.91$---and pollutes every aggregate that includes late horizons. All six of our parabola baseline arms exhibit this blow-up, which is why parabola is judged on its clean \rmood{} partition throughout this paper. The two traps point in opposite directions; neither metric family is safe alone.

\paragraph{Trap 3: zero-shot transfer is bounded by the architecture prior.}
On Physion OCP, a \emph{randomly initialized} encoder scores mean AUC $0.607$. No synthetic-training configuration exceeds it: free rollout approaches it ($0.603$), physics variants fall below it (load-bearing weighting: $0.551$), and training longer monotonically approaches---never crosses---the ceiling. Zero-shot claims must be calibrated against the random-prior floor, or they measure the architecture, not the learning.

\paragraph{Trap 4: protocol confounds fabricate results in both directions.}
Three probe-protocol choices (random vs.\ pretrained initialization, single-frame vs.\ stacked probes, probing through the projector) multiply into a false negative: the naive protocol stack reads uniform velocity decodability at $R^{2}{=}0.166$, while correcting initialization, probe window, and probe point yields $0.939$---a deficit manufactured entirely by protocol. A silent initialization bug (only $24$ of $216$ loadable encoder keys actually loaded; $192$ dropped by a parameter-renaming mismatch) invalidated an entire 45-configuration sweep before being caught by a loading guard. And converged training loss proves nothing about rollout: our clean from-scratch baseline converges to the pretrained model's prediction loss ($0.006$ vs.\ $0.005$) while its rollout OOD error is $2.8\times$ worse.

\paragraph{Checklist.}
(i)~Report per-partition \emph{and} per-horizon; never aggregate across horizons without inspecting them. (ii)~Judge with scale-aware prediction metrics (nMSE, decoded-pixel PSNR); use cosine and probe $\rho$ as diagnostics only. (iii)~Audit nMSE denominators for degenerate targets. (iv)~Calibrate transfer against a random-init control, and validate probe protocols on a random encoder.
